\subsection{Morfismos universalmente abiertos y morfismos planos.}

\begin{env}[2.4.1]
\label{IV.2.4.1-fr}
Ya hemos definido (\textbf{II}, 5.4.9) la noción de morfismo
\emph{universalmente cerrado}; del mismo modo, se dan las definiciones siguientes:
\end{env}

\begin{definition}[2.4.2]
\label{IV.2.4.2-fr}
Se dice que un morfismo de preesquemas $f:X\to Y$ es universalmente
abierto \textup{(resp.} universalmente bicontinuo, \textup{resp.} un
homeomorfismo universal) si, para todo morfismo $g:Y'\to Y$, el
morfismo $f_{(Y')}:X\times_Y Y'\to Y'$ es abierto \textup{(resp.} un homeomorfismo
sobre su imagen, \textup{resp.} un homeomorfismo sobre $Y'$).
\end{definition}

Veremos más adelante (14.3.2) que, cuando $Y$ es localmente noetheriano, la
definición de morfismo universalmente abierto dada aquí es equivalente a la
definición (\textbf{III}, 4.3.9) para los morfismos de tipo finito; el lector podrá
comprobar que no utilizamos esta última definición antes del \S~14.

\begin{proposition}[2.4.3]
\EGAReaderTarget{ega4.c04.prop.2-4-3}{IV.2.4.3-fr}
\begin{enumerate}
\item[\textup{(i)}] Una inmersión \textup{(resp.} una inmersión abierta, \textup{resp.} una inmersión
cerrada) es universalmente bicontinua \textup{(resp.} universalmente abierta, \textup{resp.}
universalmente cerrada).
\item[\textup{(ii)}] El compuesto de dos morfismos universalmente abiertos \textup{(resp.}
universalmente cerrados, \textup{resp.} universalmente bicontinuos, \textup{resp.} dos
homeomorfismos universales) también lo es.
\item[\textup{(iii)}] Si $f:X\to Y$ es un $S$-morfismo universalmente abierto \textup{(resp.}
universalmente cerrado, \textup{resp.} universalmente bicontinuo, \textup{resp.} un
homeomorfismo universal), lo mismo ocurre con
$f_{(S')}:X_{(S')}\to Y_{(S')}$ para todo cambio de base $S'\to S$.
\item[\textup{(iv)}] Si $f:X\to X'$, $g:Y\to Y'$ son dos $S$-morfismos
universalmente abiertos \textup{(resp.} universalmente cerrados, \textup{resp.} universalmente
bicontinuos, \textup{resp.} dos homeomorfismos universales), lo mismo ocurre con
$f\times_S g:X\times_S Y\to X'\times_S Y'$.
\item[\textup{(v)}] Sean $f:X\to Y$, $g:Y\to Z$ dos morfismos tales que $f$ sea
sobreyectivo; si $g\circ f$ es universalmente abierto \textup{(resp.} universalmente
cerrado, \textup{resp.} universalmente bicontinuo, \textup{resp.} un homeomorfismo universal),
lo mismo ocurre con $g$.
\item[\textup{(vi)}] Para que $f$ sea un morfismo universalmente abierto \textup{(resp.}
universalmente cerrado, \textup{resp.} universalmente bicontinuo, \textup{resp.} un
homeomorfismo universal), es necesario y suficiente que $f_{\mathrm{red}}$ lo
sea.
\item[\textup{(vii)}] Sea $(U_\alpha)$ un recubrimiento abierto de $Y$. Para que
$f:X\to Y$ sea universalmente abierto \textup{(resp.} universalmente cerrado, \textup{resp.}
universalmente bicontinuo, \textup{resp.} un homeomorfismo universal), es necesario y
suficiente que, para todo $\alpha$, su restricción
$f^{-1}(U_\alpha)\to U_\alpha$ lo sea.
\end{enumerate}
\end{proposition}

La afirmación (i) se deduce de (\textbf I, 4.3.2). La afirmación (ii) se deduce
inmediatamente de las definiciones, y lo mismo ocurre con (iii), reduciéndose
al caso en que $Y=S$, $Y'=S'$, lo que puede hacerse gracias a (\textbf I, 3.3.11);
se sabe que (iv) se deduce de (ii) y (iii) (\textbf I, 3.5.1). Para demostrar (v),
observemos que, para todo morfismo $Z'\to Z$,
$f_{(Z')}:X_{(Z')}\to Y_{(Z')}$ es sobreyectivo
(\textbf I, 3.5.2); es posible, pues, limitarse a demostrar que, si $g\circ f$ es
abierto (resp. cerrado, resp. un
\oldpage[IV-2]{20}
homeomorfismo sobre su imagen, resp. un homeomorfismo biyectivo), lo mismo
ocurre con $g$, y se trata, por tanto, de una cuestión puramente topológica.
En el caso en que $g\circ f$ es abierto (resp. cerrado), el hecho de que $g$
sea entonces abierto (resp. cerrado) se deduce de Bourbaki,
\emph{Top. gén.}, cap.~I, 3.\textsuperscript{a} ed., \S~5, núm.~1, prop.~1;
para los otros dos casos, es posible limitarse a suponer que
$g(f(X))=g(Y)=Z$, es decir, al caso en que $g\circ f$ es un homeomorfismo de
$X$ sobre $Z$; como $f$ es sobreyectivo, $g$ es necesariamente biyectivo y,
como $g$ es una aplicación continua y abierta por lo precedente, $g$ es
efectivamente un homeomorfismo de $Y$ sobre $Z$.

Para demostrar (vi), observemos que decir que un morfismo $g$ es abierto
(resp. cerrado, resp. un homeomorfismo sobre su imagen, resp. un
homeomorfismo biyectivo) equivale a decir que $g_{\mathrm{red}}$ posee la
misma propiedad. Por otra parte (\textbf I, 5.1.8), para todo morfismo $Y'\to Y$, se
tiene
$\bigl(X_{\mathrm{red}}\times_{Y_{\mathrm{red}}}Y'_{\mathrm{red}}\bigr)_{\mathrm{red}}
=(X\times_Y Y')_{\mathrm{red}}$,
por lo que la observación precedente muestra que, si $f_{\mathrm{red}}$ es
universalmente abierto (resp. universalmente cerrado, resp. universalmente
bicontinuo, resp. un homeomorfismo universal), lo mismo ocurre con $f$. El
recíproco se demuestra del mismo modo, observando aquí que, para todo morfismo
$Y''\to Y_{\mathrm{red}}$, se tiene
$\bigl(X_{\mathrm{red}}\times_{Y_{\mathrm{red}}}Y''\bigr)_{\mathrm{red}}
=(X\times_Y Y'')_{\mathrm{red}}$
(\textbf{I}, 5.1.8).

Por último, la necesidad de (vii) se deduce inmediatamente de (iii).
Recíprocamente, supongamos satisfecha la condición (vii), y sea
$g:Y'\to Y$ un morfismo; entonces los
$g^{-1}(U_\alpha)=U'_\alpha$ forman un recubrimiento abierto de $Y'$ y, si
se denota por $f_\alpha$ la restricción
$f^{-1}(U_\alpha)\to U_\alpha$ de $f$, y por $f'$ el morfismo
$f_{(Y')}$, la restricción
$f'^{-1}(U'_\alpha)\to U'_\alpha$ no es otra que
$(f_\alpha)_{(U'_\alpha)}$. Es posible, pues, limitarse a demostrar que $f$ es
abierto (resp. cerrado, resp. un homeomorfismo sobre su imagen, resp. un
homeomorfismo sobre $Y$), lo que es inmediato.
